\chapter{Топология хиггсовского вакуума и локализация смены ранга}
\label{ch:v6-spectral-transition-higgs-vacuum-topology-localization-gate}

\section{Постановка после разделения секторов}

Архитектурное решение предыдущей главы запретило использовать семейный
вихрь $(Q,T,B)$ как скрытый локализатор хиггс-разрешённой опоры. Теперь
проверяется более минимальная возможность: способен ли один стандартный
комплексный дублет Хиггса сам поддерживать топологически защищённую область
$H=0$ и тем самым локализовать переход
\begin{equation}
 \rank W_\nu(H):0\longrightarrow1,
 \qquad
 W_\nu(H)=\widetilde H\widetilde H^\dagger.
\end{equation}

Тест относится только к топологической обязательности нуля. Он не
запрещает нетопологические седла, сфалероны или динамически метастабильные
электрослабые струны.

\section{Точное вакуумное многообразие}

Для стандартного потенциала
\begin{equation}
 V(H)=\lambda\left(H^\dagger H-\frac{v^2}{2}\right)^2,
 \qquad H\in\mathbb C^2,
\end{equation}
множество минимумов равно
\begin{equation}
 \mathcal M_H=\left\{H\in\mathbb C^2:
 H^\dagger H=\frac{v^2}{2}\right\}\simeq S^3.
 \label{eq:v6-higgs-vacuum-s3}
\end{equation}
То же пространство получается как электрослабая орбита
\begin{equation}
 \mathcal M_H\simeq
 \frac{SU(2)_L\times U(1)_Y}{U(1)_{\rm em}}
 \simeq S^3
\end{equation}
для стандартного действия на одном дублете. Глобальная топология
электрослабых вакуумов и возможность иных, но уже модифицированных
глобальных форм подробно разобраны в
\cite{higgs-topology-gripaios-randal-williams2017}.

Отсюда
\begin{equation}
 \pi_0(\mathcal M_H)=0,
 \qquad
 \pi_1(\mathcal M_H)=0,
 \qquad
 \pi_2(\mathcal M_H)=0,
 \qquad
 \pi_3(\mathcal M_H)=\mathbb Z.
 \label{eq:v6-higgs-homotopy-ledger}
\end{equation}
Первые три группы классифицируют соответственно доменные стены,
линейные дефекты и точечные ядра в трёх пространственных измерениях
\cite{higgs-topology-mermin1979,higgs-topology-achucarro-vachaspati2000}.

\section{Явное стягивание кандидата струны}

На окружности вокруг предполагаемой линии рассмотрим обычный фазовый
кандидат с winding $n$:
\begin{equation}
 h_0(\varphi)=
 \begin{pmatrix}0\\ e^{in\varphi}\end{pmatrix}.
\end{equation}
Внутри $S^3$ он стягивается семейством
\begin{equation}
 h_t(\varphi)=
 \begin{pmatrix}
 \sin(\pi t/2)\\
 \cos(\pi t/2)e^{in\varphi}
 \end{pmatrix},
 \qquad 0\leq t\leq1.
 \label{eq:v6-higgs-loop-contraction}
\end{equation}
Для всех $t,\varphi$ выполнено $h_t^\dagger h_t=1$, а при $t=1$
получается постоянный вектор $(1,0)^T$. Следовательно, даже произвольная
фазовая намотка одного компонента снимается через второй компонент без
прохождения через $H=0$.

Это является точной причиной, почему вложенная электрослабая $Z$-струна
не защищена одной группой $\pi_1(\mathcal M_H)$. Такие решения могут
существовать и исследоваться динамически, но являются нетопологическими
\cite{higgs-topology-achucarro-vachaspati2000}.

\section{Явное стягивание кандидата монополя}

Пусть $n\in S^2$ и вложим сферу в вещественные координаты $S^3$ как
$(n,0)$. Гомотопия
\begin{equation}
 F_t(n)=\left(\sqrt{1-t^2}\,n,t\right),
 \qquad0\leq t\leq1,
 \label{eq:v6-higgs-sphere-contraction}
\end{equation}
остаётся на единичной трёхмерной сфере и при $t=1$ становится постоянной.
Тем самым один дублет не принуждает точечное ядро $H=0$ через
$\pi_2$. Электрослабые монопольные и струновые конфигурации, обсуждаемые
в литературе, требуют более тонкого gauge- или динамического чтения и не
изменяют этот гомотопический результат
\cite{higgs-topology-patel-vachaspati2022}.

\section{Почему текстура не создаёт переход ранга}

Ненулевая группа $\pi_3(S^3)=\mathbb Z$ допускает карты
\begin{equation}
 H:S^3_{\rm space}\longrightarrow S^3_{\rm vac}
\end{equation}
ненулевой степени. Но по определению такая карта остаётся на
$H^\dagger H=v^2/2$ во всём пространстве. Поэтому
\begin{equation}
 \rank W_\nu(H)=1
\end{equation}
в каждой точке: текстура может нести глобальное winding-чтение, но не
локализует поверхность или ядро смены ранга $0\to1$.

Кроме того, в чистой двухпроизводной скалярной энергии трёхмерная текстура
не получает фиксированного размера по балансу Деррика. Высшие производные,
gauge-динамика или дополнительный заряд могли бы изменить стабильность,
но это уже отдельный механизм и не является топологическим выводом нуля
$H$ из одного дублета.

\begin{theorem}[Отсутствие внутренней топологической локализации]
В стандартной архитектуре одного электрослабого дублета вакуумное
многообразие равно $S^3$. Его группы $\pi_0$, $\pi_1$ и $\pi_2$
тривиальны, поэтому ни стена, ни струна, ни точечный дефект не заставляют
поле проходить через $H=0$. Класс $\pi_3$ допускает текстуру, но сохраняет
ранг $W_\nu(H)$ равным единице повсюду. Следовательно, топология одного
Хиггса не локализует спектральный переход ранга $0\to1$.
\end{theorem}

\section{Граница глобальной формы группы}

Локально изоморфные электрослабые группы могут иметь другие глобальные
вакуумные многообразия, включая линзовые пространства с нетривиальной
$\pi_1$ \cite{higgs-topology-gripaios-randal-williams2017}. Но их выбор изменяет
глобальную калибровочную архитектуру и ограничения на фермионные массы.
Текущий проект использует стандартные представления
$SU(3)_c\times SU(2)_L\times U(1)_Y$ и не вывел альтернативный quotient,
который создавал бы нужный класс. Такая возможность относится к новой
модели, а не к скрытому сектору Тома VI.

\begin{nogocriterion}[Нетопологическое решение не является топологическим выводом]
Численное нахождение вложенной $Z$-струны, сфалерона или локального нуля
$H$ не закрывает настоящий гейт, если конфигурация не защищена классом
текущего вакуумного многообразия либо её устойчивость требует нового
коэффициента или поля.
\end{nogocriterion}

\section{Следующий различающий тест}

После закрытия топологической локализации остаётся только динамический
маршрут внутри спектральной ветви. Следующий гейт
\path{version6_spectral_transition_higgs_zero_finite_energy_saddle_gate}
должен проверить, допускает ли уже выведенное хиггсовское действие
конечную стационарную конфигурацию с локальным $H=0$ без топологической
защиты и без портала к $(Q,T,B)$.

Нулевая гипотеза: для одного дублета без дополнительного заряда или
граничного условия всякая конечная локальная депрессия $|H|$ либо
релаксирует к однородному вакууму, либо является седлом с отрицательной
модой и потому не задаёт устойчивую частицу.

Стоп-критерий: если масштабирование и полный радиальный гессиан не дают
локального минимума с $H(0)=0$ и $H(\infty)=H_{\rm vac}$, внутренняя
хиггсовская локализация закрывается. Тогда новая материя потребует нового
сохраняемого заряда, поля или явно расширенной конечной архитектуры.

\begin{protocol}
\begin{enumerate}
 \item вакуум одного стандартного дублета: \textbf{$S^3$};
 \item $\pi_0$, $\pi_1$, $\pi_2$: \textbf{тривиальны};
 \item фазовая петля: \textbf{явно стягивается без $H=0$};
 \item сферический кандидат: \textbf{явно стягивается без $H=0$};
 \item $\pi_3$: \textbf{$\mathbb Z$, но ранг $W_\nu$ остаётся один};
 \item топологическая локализация $0\to1$: \textbf{закрыта};
 \item нетопологические электрослабые конфигурации:
 \textbf{не исключены, но не являются выводом частицы};
 \item следующий тест: \textbf{конечная нетопологическая седловая
 конфигурация одного Хиггса};
 \item физическое замыкание: \textbf{не достигнуто}.
\end{enumerate}
\end{protocol}

Машинный сертификат сохранён в
\path{s2t_v6_spectral_transition_higgs_vacuum_topology_localization_gate_results.json}.

\begin{thebibliography}{9}
\bibitem{higgs-topology-gripaios-randal-williams2017}
B.~Gripaios, O.~Randal-Williams,
\emph{Topology of the Electroweak Vacua},
Physics Letters B 770 (2017), 309--313; arXiv:1610.05623.

\bibitem{higgs-topology-mermin1979}
N.~D.~Mermin,
\emph{The Topological Theory of Defects in Ordered Media},
Reviews of Modern Physics 51 (1979), 591--648.

\bibitem{higgs-topology-achucarro-vachaspati2000}
A.~Achucarro, T.~Vachaspati,
\emph{Semilocal and Electroweak Strings},
Physics Reports 327 (2000), 347--426; arXiv:hep-ph/9904229.

\bibitem{higgs-topology-patel-vachaspati2022}
T.~Patel, T.~Vachaspati,
\emph{Kibble Mechanism for Electroweak Magnetic Monopoles and Magnetic
Fields},
Physical Review D 106 (2022), 035041; arXiv:2108.05357.
\end{thebibliography}